\section{Quasi-coherent modules, vector bundles, and finite \v{C}ech complexes}

\subsection{Affine module data}

On an affine chart $U_i=\Spec(A_i)$, a quasi-coherent module is represented by an $A_i$-module
$M_i$.  For a finitely presented theory one usually takes a finite presentation matrix, although the
principal-cover gluing theorem only needs a certified module interface.  Restriction to a principal
open $D(f)$ is the localization $(M_i)_f$.

For a strict atlas, descent data for a quasi-coherent module consists of certified isomorphisms
\[
   \theta_{ij}:(M_i)_{f_{ij}}
      \xrightarrow{\sim}
      (M_j)_{f_{ji}}
\]
semilinear with respect to $\phi_{ij}$, together with inverse and triple-cocycle witnesses.  With
finite principal-overlap atlases, the same data are supplied on every overlap piece and checked on
finite refinements.

\begin{definition}[Certified quasi-coherent atlas module]
A certified quasi-coherent atlas module is a finite affine atlas, one certified module on each chart,
and finite localization isomorphism/cocycle certificates on the overlaps.
\end{definition}

This definition is purely finite.  Its external interpretation is the conventional quasi-coherent
sheaf obtained by descent.

\subsection{Sections and the sheaf equalizer}

A global section code is a tuple $(s_i)_{i<r}$ with $s_i\in M_i$ and overlap witnesses
\[
                    \theta_{ij}(s_i)=s_j.
\]
If a chart or overlap is represented by several principal pieces, compatibility on those pieces is
compressed by Theorem~\ref{thm:principal-gluing}.  Equality of global sections is componentwise
certificate equality; refinement invariance follows from Proposition~\ref{prop:refinement-invariance}.

For a principal cover of a single affine chart, Corollary~\ref{cor:equalizer} proves the sheaf
condition directly.  Hence the structure presheaf $D(f)\mapsto A_f$ and module presheaf
$D(f)\mapsto M_f$ are sheaves on the finite standard-open basis in exactly the finite sense used by
an iPRA proof.

\subsection{Vector bundles and transition matrices}

A rank-$d$ vector bundle on a finite atlas is represented by a free module $A_i^d$ on each chart and
invertible matrices
\[
                     G_{ij}\in\operatorname{GL}_d((A_i)_{f_{ij}})
\]
with supplied inverse matrices and cocycle witnesses $G_{ik}=G_{jk}G_{ij}$ on triple overlaps.
All matrix identities are finite coefficient identities in localizations.  A section is a tuple of
vectors $v_i\in A_i^d$ satisfying $v_j=G_{ij}v_i$.

Line bundles are the case $d=1$: transition matrices are units.  The twists $\cO(d)$ on projective
space have transition units given by integer powers of coordinate ratios.  Negative powers live in
the overlap localization and require no extra total inversion operation outside its certified
domain.

\subsection{Coherent modules and finite resolutions}

A coherent module may be represented chartwise by a finite presentation
\[
                       A_i^{u_i}\xrightarrow{R_i}A_i^{v_i}\to M_i\to0.
\]
Overlap isomorphisms are matrices on generators together with witnesses that they respect relations,
plus inverse witnesses.  Kernels, cokernels, tensor products, and duals are represented by explicit
finite matrices and relation certificates.

A finite locally free resolution is especially useful: every differential is a matrix, the equation
$d_{p-1}d_p=0$ is coefficientwise, and exactness at a supplied vector is certified by a preimage.
The framework can verify a finite resolution without deciding module equality globally.

\subsection{Finite \v{C}ech complexes}

Let $U_0,\ldots,U_{r-1}$ be a finite atlas or a finite principal cover.  For a module datum
$\mathcal F$, define the ordered \v{C}ech group in degree $p$ by the finite direct sum
\[
 \check C^p(U,\mathcal F)
   =\bigoplus_{i_0<\cdots<i_p}
     \mathcal F(U_{i_0}\cap\cdots\cap U_{i_p}).
\]
Each multiple intersection is represented by an iterated localization or a finite principal
refinement.  The differential is
\[
   (ds)_{i_0\ldots i_{p+1}}
      =\sum_{a=0}^{p+1}(-1)^a
       s_{i_0\ldots\widehat{i_a}\ldots i_{p+1}}\big|_{U_{i_0\ldots i_{p+1}}}.
\]
All sums and restrictions are fixed PR operations on finite tuples.

The identity $d^2=0$ is proved by pairing the two ways of deleting two indices; their signs are
opposite.  An iPRA proof enumerates the bounded pairs $(a,b)$ and carries a quantifier-free
coefficient accumulator.  It does not use induction on a statement quantifying over sheaves.

\subsection{Cohomology by explicit normal forms}

For a finite complex, a cohomology class may be represented in several certificate-relative ways:

\begin{enumerate}[label=\arabic*.]
\item a cocycle $z$ with a witness $dz=0$, modulo proof-relevant equality
      $z-z'=dy$;
\item a supplied direct-sum normal form with constructors for decomposition and uniqueness;
\item a finite presentation matrix for $\ker d/\im d$;
\item a contracting homotopy $h$ satisfying $hd+dh=1$ on a designated summand.
\end{enumerate}

The projective-space computation uses the fourth and second forms.  Every monomial-weight summand is
finite and has either a one-dimensional survivor or an explicit insertion homotopy.

\subsection{Relation to conventional sheaf cohomology}

The arithmetic construction gives exact \v{C}ech cohomology of a chosen finite atlas.  To identify it
with derived-functor sheaf cohomology one needs a comparison theorem, usually supplied by affine
acyclicity for quasi-coherent sheaves and acyclicity of finite intersections.  For the standard cover
of projective space, this is the standard conventional argument.

There are two honest formalization targets:

\begin{description}[style=nextline]
\item[Immediate target.]
Formalize the finite \v{C}ech complex and its normal forms in iPRA.  This already proves concrete
local-to-global statements and computes a canonical finite algebraic invariant.

\item[Extended target.]
Develop a coded homological-algebra library and an affine-acyclicity theorem, then prove that the
finite \v{C}ech result represents the usual sheaf cohomology.  This is a larger project and is not
silently assumed to be complete here.
\end{description}

The report states conventional cohomology conclusions only with this comparison boundary made
explicit.
